\chapter{Introducción y Objetivos}
\label{sec.cap1:introduccion}

La prevención y protección contra incendios en instalaciones industriales y de ingeniería electrónica constituye una
disciplina primordial orientada a salvaguardar la vida humana, preservar la integridad psicofísica de los trabajadores y
garantizar la continuidad operativa de los procesos productivos. En entornos tecnológicos donde coexisten elevadas
densidades de potencia eléctrica, conmutaciones continuas y materiales poliméricos en aislaciones y canalizaciones, el
fuego representa un riesgo permanente que demanda criterios rigurosos de ingeniería y gestión.

El presente trabajo práctico final aborda de forma integral el fenómeno del incendio desde sus fundamentos térmicos y
cinéticos, la física de las fallas eléctricas de origen común en tableros e instalaciones, y el impacto toxicológico y
fisiológico de los gases de combustión sobre las personas. Asimismo, se analizan exhaustivamente los sistemas de
detección y extinción disponibles, las alternativas de solución de ingeniería y organizativas para la prevención de
accidentes, y los requerimientos del marco regulatorio argentino fijados por la Ley Nº 19.587, el Decreto Reglamentario
351/79 y las normas IRAM y AEA.

\section{Objetivos del Informe}

Los objetivos específicos planteados para este informe técnico comprenden:
\begin{itemize}
    \item Explicar los fundamentos de la combustión mediante los modelos del triángulo y tetraedro del fuego,
          analizando los tipos de combustión, los límites de inflamabilidad y las temperaturas de ignición.
    \item Clasificar en profundidad los tipos de fuego normalizados (Clases A, B, C, D y K) conforme a la norma argentina
          IRAM 3517-1, evaluando los mecanismos de reacción, riesgos asociados y agentes extintores correspondientes.
    \item Analizar la física de las fallas de origen eléctrico que generan focos ígneos, incluyendo el efecto Joule, la
          resistencia de constricción en contactos degradados y los fenómenos de arco eléctrico en serie y en paralelo.
    \item Evaluar el impacto del humo y de gases tóxicos asfixiantes (monóxido de carbono y cianuro de hidrógeno) sobre
          el cuerpo humano y las condiciones psicofísicas de evacuación.
    \item Desarrollar casos prácticos de aplicación y cálculo reglamentario conforme al Decreto Reglamentario 351/79 (Carga
          de Fuego y Unidades de Ancho de Salida) y las prescripciones térmicas de la reglamentación AEA 90364.
    \item Analizar en detalle las tecnologías de sistemas de detección de incendios (ópticos, iónicos, de haz
          proyectado, por aspiración de alta sensibilidad, térmicos, por cable sensor lineal, de llama y de gases).
    \item Caracterizar los tipos de extintores portátiles (agentes, sistemas de presurización y potencial extintor) y los
          sistemas fijos de extinción (redes de hidrantes, rociadores automáticos, agentes limpios, agua nebulizada y
          espumas).
    \item Presentar y describir alternativas de solución de ingeniería, prevención técnica y protección para evitar
          accidentes e incendios, estructuradas bajo la jerarquía de control de riesgos y la organización humana de
          respuesta (E.P.I. y E.S.I.).
    \item Analizar con enfoque forense siniestros reales documentados (depósito Iron Mountain, subestación Edesur
          Caballito y fábrica de semiconductores Renesas Electronics), integrando las causas físicas, los cálculos
          legales y los modos de falla en las barreras de detección, extinción y gestión humana.

\end{itemize}
